%% =====================================================================
%% Tagged PDF (PDF/UA): support for packages that LaTeX's tagging code does not handle by itself yet.
%% Read by packages.tex; does something only when tagging is on (\DocumentMetadata{..., tagging=on}),
%% so the untagged editions are not affected.
%%
%% Code listings (listings): set as a code block (tag Code lines in a verbatim structure), with the
%% line numbers as labels, and the caption as the listing's caption. Adapted from the LaTeX team's own
%% set-up for the tagpdf documentation (tagpdfdocu-patches.sty, 2026-08-21, LaTeX Project Public
%% License); without it, every listing closed the section around it in the tagged structure.
%%
%% Numbered equations (MathML): multline gets a structure for its equation number, as equation
%% does, and an equation with split keeps its number in the MathML (see below and tagging.lua).
%% =====================================================================
\ExplSyntaxOn
\tag_if_active:T
  {
    \NewTaggingSocket{listings/number}{2}
    \NewTaggingSocketPlug{listings/number}{default}
      {
        \tag_mc_end_push:
        \tagstructbegin{tag=Lbl}
        \tagmcbegin{}
        #2
        \tagmcend
        \tagstructend
        \tag_mc_begin_pop:n{}
      }
    \AssignTaggingSocketPlug{listings/number}{default}
    \DeclareInstance{blockenv}{lstlisting}{std}
      {
        name                   = lstlisting,
        tag-name               = verbatim,
        tag-attr-class         = ,
        tagging-recipe         = standard,
        inner-level-counter    = ,
        transparent-level      = true,
        block-instance         = lstlisting,
        inner-instance         = ,
        final-code             = \AssignStructureRole{para/textblock}
                                   {\UseStructureName{block/verbatim/codeline}}\ignorespaces,
        tagging-suppress-paras = true
      }
    \DeclareInstance{block}{lstlisting-1}{std}
      {
        left-margin = 0pt, para-vspace = 0pt, begin-vspace = 0pt, end-vspace = 0pt,
        begin-extra-vspace = 0pt, end-extra-vspace = 0pt, begin-unchained-vspace = 0pt,
        % listings adds its own skips and break points (\lstset aboveskip, belowskip)
        begin-penalty = 10000, end-penalty = 10000
      }
    \DeclareInstanceCopy{block}{lstlisting-2}{lstlisting-1}
    \DeclareInstanceCopy{block}{lstlisting-3}{lstlisting-1}
    \DeclareInstanceCopy{block}{lstlisting-4}{lstlisting-1}
    \DeclareInstanceCopy{block}{lstlisting-5}{lstlisting-1}
    \DeclareInstanceCopy{block}{lstlisting-6}{lstlisting-1}
    \tagpdfsetup{float/new=lstlisting}
    % multline: luamml leaves multline out of the environments whose equation number gets its
    % own structure, but still refers to that structure when it writes the MathML, which gives
    % "structure with label 0 is unknown". Treat multline like equation and gather.
    % Both only when the MathML is made (luamml, math/setup in \DocumentMetadata).
    \AddToHook{begindocument}
      {
        \IfPackageLoadedT{luamml}
          {
            \tl_if_exist:cF { c__luamml_label_multline_tl }
              { \tl_const:cn { c__luamml_label_multline_tl } { } }
            % equation with split: see tagging.lua.
            \directlua{dofile(kpse.find_file('_options/tagging.lua'))}
            \NewTaggingSocketPlug{math/luamml/mtable/tag/save}{thesis}
              {
                \bool_if:NTF \l__math_gathersplit_bool
                  { \use:c{thesis@luamml@savetag} \l__luamml_amsmath_tag_struct_num_int }
                  { \__luamml_amsmath_save_tag_with_struct_elem:N \l__luamml_amsmath_tag_struct_num_int }
              }
            \AssignTaggingSocketPlug{math/luamml/mtable/tag/save}{thesis}
            \NewTaggingSocketPlug{math/luamml/mtable/tag/set}{thesis}
              { \bool_if:NF \l__math_gathersplit_bool { \__luamml_amsmath_set_tag: } }
            \AssignTaggingSocketPlug{math/luamml/mtable/tag/set}{thesis}
            \NewTaggingSocketPlug{math/luamml/mtable/innertable/finalize}{thesis}
              {
                \__luamml_amsmath_finalize_inner_table:
                \bool_if:NT \l__math_gathersplit_bool { \use:c{thesis@luamml@wraptag} }
              }
            \AssignTaggingSocketPlug{math/luamml/mtable/innertable/finalize}{thesis}
          }
      }
    % Theorems, proofs and other blocks that stand on their own: LaTeX's tagging code ends them
    % without the usual note that what follows straight on continues the paragraph (\if@endpe),
    % which the tagging needs, but it changes the look: text straight after was indented, and a
    % theorem, proof or list straight after got \partopsep more space above it. Here the look
    % stays as in the untagged editions, while the tagging code still sees what it expects.
    % (LaTeX's code as of June 2026, with the socket block/endpe; the November 2026 code is handled below)
    \socket_if_exist:nT { block/endpe }
      {
    \bool_new:N \g__thesis_block_endpe_bool   % a block on its own has just ended
    \bool_new:N \l__thesis_block_pre_bool
    \bool_new:N \l__thesis_block_space_bool
    \bool_new:N \l__thesis_block_after_bool
    \cs_new_protected:Npn \__thesis_block_after_standalone:
      {
        \bool_gset_true:N \g__thesis_block_endpe_bool
        \def \par
          {
            \use:c { @restorepar } \bool_gset_false:N \g__thesis_block_endpe_bool
            \everypar { } \par
          }
        \everypar
          {
            { \setbox 0 \lastbox }
            \everypar { } \bool_gset_false:N \g__thesis_block_endpe_bool
          }
      }
    \socket_new_plug:nnn { block/endpe } { thesis-standalone }
      { \legacy_if_gset_false:n { @endpe } \aftergroup \__thesis_block_after_standalone: }
    \AddToHook{begindocument}
      {
        \cs_new_eq:NN \__thesis_block_recipe_standalone: \__block_recipe_standalone:
        \cs_gset_protected:Npn \__block_recipe_standalone:
          {
            \__thesis_block_recipe_standalone:
            \socket_assign_plug:nn { block/endpe } { thesis-standalone }
          }
        \NewTaggingSocketPlug { block/recipe } { thesis }
          {
            \legacy_if:nTF { @endpe }
              { \bool_set_true:N \l__thesis_block_pre_bool }
              { \bool_set_false:N \l__thesis_block_pre_bool }
            \bool_set_eq:NN \l__thesis_block_space_bool \l__thesis_block_pre_bool
            \bool_if:NT \g__thesis_block_endpe_bool
              { \bool_set_true:N \l__thesis_block_space_bool }
            \bool_gset_false:N \g__thesis_block_endpe_bool
            \use:c { __block_recipe_#1: }
            \legacy_if:nTF { @endpe }
              { \bool_set_true:N \l__thesis_block_after_bool }
              { \bool_set_false:N \l__thesis_block_after_bool }
            \bool_if:NT \l__thesis_block_space_bool
              {
                % the space above as without tagging, then the state the tagging code set
                \legacy_if_gset_true:n { @endpe }
                \cs_set_eq:NN \__thesis_displayblock_begin: \__kernel_displayblock_begin:
                \cs_set_protected:Npn \__kernel_displayblock_begin:
                  {
                    \bool_if:NTF \l__thesis_block_after_bool
                      { \legacy_if_gset_true:n { @endpe } }
                      { \legacy_if_gset_false:n { @endpe } }
                    \__thesis_displayblock_begin:
                  }
              }
          }
        \AssignTaggingSocketPlug { block/recipe } { thesis }
      }
      }
    % Figures side by side (a minipage each, with its own caption, in one figure): a figure may
    % have only one caption in a tagged PDF, so each minipage after the first starts a figure of its
    % own in the structure (the pages are unchanged).
    \int_new:N \g__thesis_float_captions_int
    \AddToHook{cmd/@floatboxreset/after}{ \int_gzero:N \g__thesis_float_captions_int }
    \AddToHook{cmd/end@float/before}{ \int_gzero:N \g__thesis_float_captions_int }
    \AddToHook{cmd/end@dblfloat/before}{ \int_gzero:N \g__thesis_float_captions_int }
    \AddToHook{cmd/caption/before}{ \int_gincr:N \g__thesis_float_captions_int }
    \AddToHook{env/minipage/before}
      {
        \bool_lazy_and:nnT
          { \cs_if_exist_p:c { @captype } }
          { \int_compare_p:nNn { \g__thesis_float_captions_int } > { 0 } }
          {
            \tagpdfsetup{float/split}
            \int_gzero:N \g__thesis_float_captions_int
          }
      }
    \cs_gset:Npn \thesistagging {}
  }
\ExplSyntaxOff
\ifdefined\thesistagging
\makeatletter
\lst@Key{numbers}{none}{%
    \let\lst@PlaceNumber\@empty
    \lstKV@SwitchCases{#1}%
    {none:\\%
     left:\def\lst@PlaceNumber{\llap{\normalfont
                \UseTaggingSocket{listings/number}{}{\lst@numberstyle{\thelstnumber}}\kern\lst@numbersep}}\\%
     right:\def\lst@PlaceNumber{\rlap{\normalfont
                \kern\linewidth \kern\lst@numbersep
                \UseTaggingSocket{listings/number}{}{\lst@numberstyle{\thelstnumber}}}}%
    }{\PackageError{Listings}{value `#1' for keyword `numbers' not supported}\@ehc}}
\def\lst@EveryDisplay{%
    \MathCollectFalse
    \parskip\z@skip % no extra stretch above the listing (the block adds the paragraph skip)
    \ifx\lst@float\relax \SimpleBlockEnv{lstlisting}{}%
    \else \SimpleBlockEnv{lstlisting}{tagging-recipe = standalone}\fi
    \thesis@unskipzero} % the block's skips of nothing: they stayed at the foot of a page that breaks here
\def\thesis@unskipzero{\let\thesis@next\relax
    \ifvmode\ifnum\lastnodetype=11 \ifdim\lastskip=\z@ \ifdim\gluestretch\lastskip=\z@
    \ifdim\glueshrink\lastskip=\z@ \unskip\let\thesis@next\thesis@unskipzero\fi\fi\fi\fi\fi
    \thesis@next}
\lst@AddToHook{PreInit}{\ifhmode \UseTaggingSocket{inline/begin}{tag=\UseStructureName{text/verb}}\fi}
\lst@AddToHook{DeInit}{\ifhmode \UseTaggingSocket{inline/end}%
    \else \edef\@current@float@struct{\UseName{tag_get:n}{struct_num}}\BlockEnvEnd \fi}
\def\lst@makecaption{\AssignStructureRole{para/textblock}{Span}\def\@captype{lstlisting}\@makecaption}
\def\lst@maketitle{\def\@captype{lstlisting}%
    \AssignStructureRole{para/textblock}{Span}%
    \AssignTaggingSocketPlug{caption/label/begin}{noop}%
    \AssignTaggingSocketPlug{caption/label/end}{noop}%
    \AssignSocketPlug{caption/label}{noop}%
    \@makecaption\lst@title@dropdelim}
\makeatother
\fi
